\section{Conclusions}\label{sec:conclusions}

Let us conclude by discussing a handful of open questions.
\begin{itemize}
	\item The central open question is that of Dvo\v{r}\'ak and T\r{u}ma~\cite{DvorakT13}: Can one design a dynamic data structure for $\FO$ model-checking on any fixed  class of bounded expansion with polylogarithmic amortized update time? A natural approach to this question is to try to develop dynamic counterparts for various techniques of Sparsity underlying the known static algorithms~\cite{DvorakKT13,grohe2011methods,PilipczukST18}, such as transitive--fraternal augmentations, weak coloring numbers, or low treedepth colorings. After a few years of attempts, we consider this route hopeless. Instead, we believe that developing a new model-checking algorithm for $\FO$ on classes of bounded expansion, which would be more in the spirit of progressive exploration, could be a way to go.
	\item While the data structure of Dvo\v{r}\'ak and T\r{u}ma~\cite{DvorakT13} can be also deployed on any nowhere dense graph class, with the resulting amortized update time becoming $\Oh_{\Cc,H,\delta}(n^\delta)$ for any $\delta>0$, this is not the case for our data structures of \cref{thm:main-domset-be,thm:main-indset-be}. The reason is that in the nowhere dense setting, the fraternal augmentations discussed in \cref{sec:augmentations} have maximum outdegree $\Oh_{\Cc,r,\delta}(n^\delta)$ instead of a constant, and hence the Inclusion--Exclusion formula postulated in \cref{lem:IE} may run over $2^{\Oh_{\Cc,r,\delta}(n^\delta)}$ elements of $\cal M$, which is superpolynomial. We leave it open whether \drds{r} and \dris{r} on nowhere dense classes admit dynamic data structures with amortized update time $\Oh_{\Cc,r,k,\delta}(n^\delta)$, for any $\delta>0$. However, we remark that the \cref{lem:IE} is the only reason why this argument does not lift to nowhere dense classes and multiple other results, in particular \cref{thm:dvorak-tuma-examples}, can be lift verbatim.
	\item Can our data structures be derandomized? At this point, we crucially rely on randomization in the fingerprinting technique, to turn counting data structures into example-reporting data structures.
	\item It remains open whether the result of \cref{thm:apx-dynamic} can be extended to distance-$r$ domination for $r>1$ under the assumption that $G$ belongs at all times to a fixed class of bounded expansion $\Cc$. As discussed in \cref{sec:overview-other}, this is connected to the dynamic maintenance of the set system of weak $r$-reachability sets.
\end{itemize}
